\subsection{滤过双代数}

定义 2. 设 E 为带有余积 c 的双代数。~与 E 上双代数结构相容的滤过是指 E 的子模的递增序列 \((\mathrm{E}_{n})_{n\geq0}\)，满足

\[
\begin{array}{c} \mathrm{E} _ {0} = \mathrm{K}. 1, \qquad \mathrm{E} = \bigcup_ {n \geqslant 0} \mathrm{E} _ {n} \\ \mathrm{E} _ {m}. \mathrm{E} _ {n} \subset \mathrm{E} _ {m + n} \quad \text {for} \quad m \geqslant 0, n \geqslant 0 \\ c (\mathrm{E} _ {n}) \subset \sum_ {i + j = n} \mathrm{Im} (\mathrm{E} _ {i} \otimes \mathrm{E} _ {j}) \quad \text {for} n \geqslant 0. \dagger \end{array}\tag{6}
\]

带有与其双代数结构相容的滤过的双代数称为滤过双代数。

例. 设 E 为分次双代数（《代数》，第 III 章，第 11 节，第 4 号，定义 3），\((\mathrm{E}^n)_{n\geqslant 0}\) 为其分次。我们记 \(\mathrm{E}_n = \sum_{i = 0}^{n}\mathrm{E}^i\)。序列 \((\mathrm{E}_n)\) 是与 E 上双代数结构相容的滤过。

命题 6. 设 E 为滤过双代数，\((\mathrm{E}_n)_{n\geqslant 0}\) 为其滤过。对每个整数 \(n\geqslant 0\)，设 \(\mathrm{E}_n^+ = \mathrm{E}_n\cap \mathrm{E}^+\)。那么 \(\mathrm{E_0^+} = \{0\}\) 且

\[
c ^ {+} (\mathrm{E} _ {n} ^ {+}) \subset \sum_ {i = 1} ^ {n - 1} \operatorname{Im} (\mathrm{E} _ {i} ^ {+} \otimes \mathrm{E} _ {n - i} ^ {+}) \quad \text { for } n \geqslant 0. \dagger\tag{7}
\]

由于 \(\mathrm{E}_0 = \mathrm{K}.1, \mathrm{E}_0^+ = 0\)。若 \(x \in \mathrm{E}_n\)，则 \(\pi(x) = x - \varepsilon(x).1\)（公式 (1)），由此 \(\pi(x) \in \mathrm{E}_n^+\)，即 \(\pi(\mathrm{E}_n) \subset \mathrm{E}_n^+\)。于是 \(\pi \otimes \pi\) 把 \(\operatorname{Im}(\mathrm{E}_i \otimes \mathrm{E}_j)\) 映入 \(\operatorname{Im}(\mathrm{E}_i^+ \otimes \mathrm{E}_j^+)\)（对 \(i \geqslant 0, j \geqslant 0\)）。由于 \(c^+ = (\pi \otimes \pi) \circ c\) 在 \(\mathrm{E}^+\) 中成立（第 1 号，命题 2），由 (6) 得

\[
c ^ {+} \left(\mathrm{E} _ {n} ^ {+}\right) \subset \sum_ {i = 0} ^ {n} \operatorname{Im} \left(\mathrm{E} _ {i} ^ {+} \otimes \mathrm{E} _ {n - i} ^ {+}\right) = \sum_ {i = 1} ^ {n - 1} \operatorname{Im} \left(\mathrm{E} _ {i} ^ {+} \otimes \mathrm{E} _ {n - i} ^ {+}\right).
\]

推论. \(\mathbf{E}_1^+\) 的元素都是本原的。

若 \(x \in \mathrm{E}_{1}^{+}\)，则由 (7) 得 \(c^{+}(x) = 0\)，由此得 (5)。

